\documentclass[10pt,a4paper]{amsart}
\usepackage[margin=19mm]{geometry}
\usepackage{amsmath,amssymb,mathtools}
\usepackage[scr=rsfs,cal=euler]{mathalfa}
\usepackage{xfrac}
\usepackage[unicode,hidelinks,bookmarks=false]{hyperref}
\newcommand{\ie}{i.e.\ }
\newcommand{\cf}{cf.\ }
\newcommand{\resp}{respectively\ }
\newcommand{\Subsection}{\subsection}
\providecommand{\nobreakdash}{\nobreak\hbox{-}}
\setlength{\emergencystretch}{2em}
\multlinegap0pt
\usepackage{bm}
\usepackage{bbm}
\usepackage{dsfont}
\usepackage{xspace}
\usepackage{cancel}
\usepackage{mathtools}
\usepackage{amsthm}
\makeatletter
\@ifundefined{theorem}{\newtheorem{theorem}{Theorem}}{}
\@ifundefined{lemma}{\newtheorem{lemma}[theorem]{Lemma}}{}
\@ifundefined{proposition}{\newtheorem{proposition}[theorem]{Proposition}}{}
\makeatother
\newcommand{\R}{\mathbb R}
\newcommand{\dd}{\,d}
\newcommand{\loc}{\mathrm{loc}}
\title[Optimal regularity at the free boundary in one-dimensional f]{Optimal regularity at the free boundary in one-dimensional first-order mean field games}
\author{Sebastian Munoz}
\date{Source: arXiv:2605.06871v1 (CC BY 4.0)}
\begin{document}
\maketitle

\noindent\emph{Source and licence.} Optimal regularity at the free boundary in one-dimensional first-order mean field games. Version arXiv:2605.06871v1 (CC BY 4.0), distributed under the Creative Commons Attribution 4.0 International licence, as recorded in the arXiv licence metadata for this exact version and shown on the abstract page https://arxiv.org/abs/2605.06871v1. Full rights notice: licenses/arxiv-2605.06871-CC-BY-4.0.txt.\par\smallskip

\noindent\textit{Editorial scope.} Counted unit: the Proposition whose source label is \texttt{prop:finite-endpoint-regularity} in the article (source0.tex lines 999--1020, source label \texttt{prop:finite-endpoint-regularity}), delivered together with the article's own complete original proof of it (source0.tex lines 1022--1243, one \texttt{proof} environment of 222 lines). No mathematics is written, repaired or re-derived: statement and proof are the article's own text line for line. Required context retained uncounted: eq:mfg-system (lines 47--56); ass:mfg-terminal (lines 59--63); ass:planning-terminal (lines 65--68); ass:initial-data (lines 144--149); ass:h-second (lines 158--162); ass:planning-basic (lines 164--169); ass:terminal-data (lines 171--177); thm:higher-boundary-regularity (lines 188--219); eq:mass-relation (lines 335--338); eq:h-def (lines 388--392); eq:Z-def (lines 443--448); prop:endpoint-schauder-bootstrap (lines 782--793); eq:bootstrap-linear (lines 803--808); eq:reg-sing-ode (lines 904--907); eq:reg-sing-coefs (lines 909--914); lem:regular-singular-descent (lines 923--935). Every definition, display and label of those lines is retained unchanged. Omitted from this package and not redistributed: 1--46, 57--58, 64--64, 69--143, 150--157, 163--163, 170--170, 178--187, 220--334, 339--387, 393--442, 449--781, 794--802, 809--903, 908--908, 915--922, 936--998, 1021--1021, 1244--1703. Nothing inside a retained excerpt is omitted or altered: every retained body is delivered exactly as the article's lines are written, so no omission receipt of any kind accompanies this package. Citations: the retained excerpts carry no citation command at all, so no bibliography entry of the article is transcribed and no reference is redistributed. Reference adaptation: none was needed and none was made; every reference inside the delivered unit resolves inside it, so no unresolved reference or citation is produced. Printed slips kept literal: no printed word, symbol, hypothesis or number of the counted statement or of its proof was altered. Layout and class: the article's own class is \texttt{amsart} with packages including geometry, amsmath, amssymb, amsthm, xcolor, hyperref; the wrapper is the lane's pinned \texttt{amsart} editorial wrapper at 10pt on A4, which restates the theorem-like environments the retained text needs and adds the article's own macro definitions that the retained text needs (\texttt{\textbackslash R}, \texttt{\textbackslash dd}, \texttt{\textbackslash loc}), with extra wrapper packages (bm, bbm, dsfont, xspace, cancel, mathtools). The delivered numbering is the unit's own. Assets: no retained span contains a figure environment, a graphics-inclusion command, a PDF-page inclusion, a table or a TikZ picture; no image file is copied into this package, the package declares no asset dependency and no image file is redistributed with it. Licence: version arXiv:2605.06871v1 of this article is distributed under the Creative Commons Attribution 4.0 International licence, as recorded in the arXiv licence metadata for this exact version and shown on the abstract page https://arxiv.org/abs/2605.06871v1. A full rights notice is delivered at licenses/arxiv-2605.06871-CC-BY-4.0.txt. Characters: every retained excerpt is pure ASCII, so the delivered engine, XeLaTeX with its default Latin Modern text font, reports no missing character.
\begin{equation}
\label{eq:mfg-system}
\begin{cases}
        -u_t+\frac12 u_x^2=m^\theta
                &\text{in }\R\times(0,T),\\
        m_t-(mu_x)_x=0
                &\text{in }\R\times(0,T),\\
        m(\cdot,0)=m_0,
\end{cases}
\end{equation}

\begin{equation}
\label{ass:mfg-terminal}
        u(\cdot,T)=g(m(\cdot,T)),\qquad
        g(s)=c_1s^\theta,\qquad c_1\ge0 .
\end{equation}

\begin{equation}
\label{ass:planning-terminal}
        m(\cdot,T)=m_T .
\end{equation}

\begin{equation}
\label{ass:initial-data}
        p_0\in C([a,b])\cap C^{1,\sigma}(a,b),\qquad
        p_0>0\text{ in }(a,b),\qquad
        p_0(a)=p_0(b)=0,
\end{equation}

\begin{equation}
\label{ass:h-second}
        p_0''\ge -K_0\text{ in }(a,b),\qquad
        p_0''\le0\text{ in }(a,a+\delta)\cup(b-\delta,b),
\end{equation}

\begin{equation}
\label{ass:planning-basic}
        m_T\ge0,\qquad \{m_T>0\}=(a_1,b_1),\qquad
        p_T:=m_T^\theta\in
        C([a_1,b_1])\cap C^{1,\sigma}(a_1,b_1),
\end{equation}

\begin{equation}
\label{ass:terminal-data}
        C_0^{-1}\operatorname{dist}(x,\{a_1,b_1\})
        \le p_T(x)\le
        C_0\operatorname{dist}(x,\{a_1,b_1\}),
        \qquad x\in(a_1,b_1).
\end{equation}

\begin{theorem}[Regularity of the value function, pressure, and free boundary]
\label{thm:higher-boundary-regularity}
Let \((u,m)\) be a solution of \eqref{eq:mfg-system} with either
\eqref{ass:mfg-terminal} or \eqref{ass:planning-terminal}. Assume
\eqref{ass:initial-data}--\eqref{ass:h-second}, and also
\eqref{ass:planning-basic}, \eqref{ass:terminal-data} in the planning case.
Let \(\gamma\) be the Lagrangian flow and set
\(\gamma_L:=\gamma(a,\cdot)\), \(\gamma_R:=\gamma(b,\cdot)\), so that
\(\{m>0\}=\{(x,t):\gamma_L(t)<x<\gamma_R(t)\}\). Set \(p:=m^\theta\). Then,
for every compact interval \(J\subset(0,T)\),
\[
\begin{gathered}
        \gamma_L,\gamma_R\in C^\infty_{\loc}(0,T),\\
        p\in W^{1,\infty}_{\loc}(\R\times(0,T))
          \cap C^{1,\sigma}\bigl(\overline{\{m>0\}}\cap(\R\times J)\bigr),\\
        u\in C^{1,1/2}_{\loc}(\R\times(0,T))
          \cap C^{2,\sigma}\bigl(\overline{\{m>0\}}\cap(\R\times J)\bigr).
\end{gathered}
\]
If, in addition, \(p_0\in C^{\ell,\sigma}([a,b])\) for some integer
\(\ell\ge2\), then, for every compact interval
\(J\subset(0,T)\),
\[
        p\in
        C^{\ell,\sigma}\bigl(\overline{\{m>0\}}\cap(\R\times J)\bigr),
        \qquad
        u\in
        C^{\ell+1,\sigma}\bigl(\overline{\{m>0\}}\cap(\R\times J)\bigr).
\]
In particular, if \(p_0\) is smooth on \([a,b]\), then \(p\) and \(u\) are
smooth up to the free boundary from the positive phase.
\end{theorem}

\begin{equation}
\label{eq:mass-relation}
        p(\gamma(y,t),t)=p_0(y)\gamma_y(y,t)^{-\theta}
\end{equation}

\begin{equation}
\label{eq:h-def}
        h(y):=\frac{p_0(y)}{y}\quad(y\in(0,\delta)),\qquad
        h(0):=p_0'(0^+),
\end{equation}

\begin{equation}
\label{eq:Z-def}
        Z:=\gamma_y^{-(\theta+1)},\qquad
        \beta(Z):=\frac1{\theta+1}Z^{-1-\frac1{\theta+1}},\qquad
        c_\theta:=\frac{\theta}{\theta+1} .
\end{equation}

\begin{proposition}[Schauder bootstrap and smoothness of the axis trace]
\label{prop:endpoint-schauder-bootstrap}
Under the assumptions of Theorem~\ref{thm:higher-boundary-regularity},
let \((u,m)\) be the corresponding solution, \(\gamma\) the associated
Lagrangian flow, \(Z\) as in \eqref{eq:Z-def},
\(\widetilde Z(r,t):=Z(r^2/4,t)\), and \(I\Subset(0,T)\). Then for every
\(\tau\in(0,1)\) and every \(j\ge0\),
\[
        \partial_t^j\widetilde Z\in C^{1,\tau}_{\loc}([0,r_0)\times I).
\]
In particular, \(\widetilde Z(0,\cdot)\in C^\infty_{\loc}(I)\).
\end{proposition}

\begin{equation}
\label{eq:bootstrap-linear}
        \partial_t\!\left(W\beta(\widetilde Z)(\widetilde Z_j)_t\right)
        +\partial_r\!\left(WA(\widetilde Z_j)_r\right)
        =-Wf_j+\partial_t(W\mathfrak F_j),
\end{equation}

\begin{equation}
\label{eq:reg-sing-ode}
        y V_y+b(y)V=F,
\end{equation}

\begin{equation}
\label{eq:reg-sing-coefs}
        b(y):=\frac{(1+c_\theta)p_0'(y)}{c_\theta h(y)},
        \qquad
        F:=-\frac{\partial_t(\beta(Z)Z_t)+p_0''(y)Z}{c_\theta h(y)} .
\end{equation}

\begin{lemma}[H\"older estimate for a regular-singular ODE]
\label{lem:regular-singular-descent}
Let \(k\ge0\), \(\alpha\in(0,1)\), \(b\in C^{k,\alpha}([0,\rho])\), and
\(b\ge b_0>0\). Suppose that \(V\in L^\infty((0,\rho)\times I)\) satisfies
\eqref{eq:reg-sing-ode} distributionally in \((0,\rho)\times I\), with
\(F\in C^{k,\alpha}([0,\rho]\times I)\).
Then
\[
        V\in C^{k,\alpha}([0,\rho]\times I),
        \qquad
        y\,\partial_y^{k+1}V\in C^{0,\alpha}([0,\rho]\times I).
\]
\end{lemma}

\begin{proposition}[Higher-order Lagrangian endpoint regularity]
\label{prop:finite-endpoint-regularity}
Suppose, in addition to the assumptions of
Theorem~\ref{thm:higher-boundary-regularity}, that
\(p_0\in C^{\ell,\sigma}([0,\delta])\) for some integer \(\ell\ge2\) and
some \(\sigma\in(0,1)\). Let \((u,m)\) be the corresponding solution and
\(\gamma\) the associated Lagrangian flow, and set
\(\gamma_L:=\gamma(0,\cdot)\). Let \(Z\) be as in \eqref{eq:Z-def} and
fix \(I\Subset(0,T)\). Then
\[
        \gamma-\gamma_L,\quad
        \partial_t(\gamma-\gamma_L)\in
        C^{\ell,\sigma}_{\loc}([0,r_0^2/4)\times I),
\]
and
\[
        (y,t)\mapsto m(\gamma(y,t),t)^\theta
        \in C^{\ell,\sigma}_{\loc}([0,r_0^2/4)\times I).
\]
Moreover, \(Z(0,\cdot)\in C^\infty_{\loc}(I)\). The symmetric statement
holds at the right endpoint \(b\), after replacing \(y\) by \(b-y\).
\end{proposition}

\begin{proof}
With \(h\) as in \eqref{eq:h-def}, the higher regularity of \(p_0\) gives
\(h\in C^{\ell-1,\sigma}([0,\delta))\) with \(h(0)>0\). We claim that, for
\(0\le i\le2\ell-1\),
\begin{equation}
\label{eq:tangential-z-bound}
        \partial_t^i Z\in C^{0,1}_{\loc}([0,r_0^2/4)\times I),
\end{equation}
and \(Z(0,\cdot)\in C^\infty_{\loc}(I)\). The substitution \(r=2\sqrt y\)
is degenerate at the axis, so this does not follow from the H\"older
bounds on \(\partial_t^i\widetilde Z\) of Proposition
\ref{prop:endpoint-schauder-bootstrap} by direct composition; the missing
rate of vanishing at the axis, \((\partial_t^i\widetilde Z)_r=O(r)\), comes from integrating the
divergence-form equation for \(\partial_t^i\widetilde Z\) against the radial weight.

Set \(\widetilde Z_i:=\partial_t^i\widetilde Z\); with the conventions
\(\mathfrak F_0=0\), \(f_0=D\widetilde Z\) introduced in the proof of
Proposition~\ref{prop:endpoint-schauder-bootstrap},
equation~\eqref{eq:bootstrap-linear} holds for all \(i\ge0\). Integrate
it in \(r\) over \((0,r)\) at fixed \(t\); the boundary term at
\(\rho=0\) vanishes since \(W(\rho)\sim\rho^{N-1}\to0\) and
\((\widetilde Z_i)_r\) is bounded by
Proposition~\ref{prop:endpoint-schauder-bootstrap}, giving
\[
        W(r)A(r)(\widetilde Z_i)_r(r,t)
        =-\int_0^rW(\rho)f_i(\rho,t)\dd\rho
        +\partial_t\!\int_0^rW(\rho)\bigl[\mathfrak F_i-\beta(\widetilde Z)(\widetilde Z_i)_t\bigr](\rho,t)\dd\rho .
\]
The first integrand is \(O(\rho^{N-1})\). For the second term, differentiate
under the integral; the differentiated integrand is
\[
        W(\rho)\,\partial_t\!\left[
        \mathfrak F_i-\beta(\widetilde Z)(\widetilde Z_i)_t
        \right](\rho,t),
\]
which is also \(O(\rho^{N-1})\), since
Proposition~\ref{prop:endpoint-schauder-bootstrap} bounds
\(\partial_t^k\widetilde Z\) for \(k\le i+2\). Both right-hand terms are therefore
\(O(r^N)\). Dividing by \(W(r)A(r)\sim r^{N-1}\) gives \((\widetilde Z_i)_r(r,t)=O(r)\),
and so
\[
        \widetilde Z_i(r,t)-\widetilde Z_i(0,t)
        =\int_0^r(\widetilde Z_i)_r(\rho,t)\dd\rho
        =O(r^2) .
\]
With \(r=2\sqrt y\), this is \(O(y)\). Moreover
\begin{equation}
\label{eq:Vi-bounded}
        V_i(y,t):=\partial_y\partial_t^iZ(y,t)
        =\frac2r(\widetilde Z_i)_r(r,t)=O(1),
\end{equation}
so \(\partial_t^iZ\) is Lipschitz in~\(y\); it is Lipschitz in~\(t\) as
well, since Proposition~\ref{prop:endpoint-schauder-bootstrap} bounds
\(\partial_t^{i+1}\widetilde Z\). Hence \(\partial_t^iZ\) is jointly Lipschitz on
\([0,r_0^2/4)\times I\), establishing \eqref{eq:tangential-z-bound}. The
smoothness of \(Z(0,\cdot)\) is immediate from
\(Z(0,t)=\widetilde Z(0,t)\in C^\infty_{\loc}(I)\), again by
Proposition~\ref{prop:endpoint-schauder-bootstrap}.

For \(0\le i\le2\ell-3\), differentiating \eqref{eq:reg-sing-ode} \(i\)
times in \(t\) shows that \(V_i\) satisfies
\begin{equation}
\label{eq:reg-sing-Vi}
        y(V_i)_y+b(y)V_i=F_i,
        \qquad F_i:=\partial_t^iF,
\end{equation}
with \(b\) and \(F\) as in \eqref{eq:reg-sing-coefs}. Since
\(c_\theta\), \(h(y)\), and \(p_0''(y)\) are independent of \(t\),
\begin{equation}
\label{eq:Fi-explicit}
        F_i=-\frac{\partial_t^{i+1}(\beta(Z)Z_t)
        +p_0''(y)\,\partial_t^iZ}{c_\theta h(y)} ,
\end{equation}
which is a polynomial in \(Z\) and its tangential derivatives
\(\partial_t^kZ\) for \(k\le i+2\), times \(y\)-coefficients built from
\(1/h(y)\) and \(p_0''(y)\). Since \(p_0\in C^{\ell,\sigma}\) and
\(h\in C^{\ell-1,\sigma}\) with \(h\ge h_0>0\), these \(y\)-coefficients
lie in \(C^{\ell-2,\sigma}([0,r_0^2/4))\); the principal coefficient
\(b\) in \eqref{eq:reg-sing-coefs} is a quotient of \(p_0'\) and \(h\), both
in \(C^{\ell-1,\sigma}\), and so itself lies in
\(C^{\ell-1,\sigma}([0,r_0^2/4))\). After decreasing \(r_0\) if
necessary, \(b\ge b_0>0\) on the endpoint interval, and the binding
regularity constraint for
Lemma~\ref{lem:regular-singular-descent} comes from \(F_i\) rather than
from \(b\).
For \(n=0\), \eqref{eq:tangential-z-bound} bounds the tangential factors
of \(F_i\) in \eqref{eq:Fi-explicit} (with \(k\le i+2\le2\ell-1\)),
placing \(F_i\) in \(C^{0,\sigma}_{\loc}([0,r_0^2/4)\times I)\); and
\(V_i\) is bounded by \eqref{eq:Vi-bounded}. Lemma
\ref{lem:regular-singular-descent} therefore gives
\begin{equation}
\label{eq:Vi-base}
        V_i,\ y(V_i)_y\in C^{0,\sigma}_{\loc}([0,r_0^2/4)\times I),
        \qquad i\le2\ell-3 .
\end{equation}
For higher normal derivatives, the strategy is to apply Lemma
\ref{lem:regular-singular-descent} to \eqref{eq:reg-sing-Vi} at order
\(n\), which requires \(F_i\in C^{n,\sigma}_{\loc}([0,r_0^2/4)\times I)\).
We prove the resulting estimate
\begin{equation}
\label{eq:Vi-n-form}
        \partial_y^nV_i,\ y\,\partial_y^{n+1}V_i
        \in C^{0,\sigma}_{\loc}([0,r_0^2/4)\times I),
        \qquad i+2(n+1)\le2\ell-1,
\end{equation}
by induction on \(n\). The 2:1 trade-off in the constraint reflects the
structure of \(F\) in \eqref{eq:reg-sing-coefs}: the term
\(\partial_t(\beta(Z)Z_t)\) makes \(F_i\) depend on tangential derivatives
of \(Z\) up to order \(i+2\), and each additional level of normal-descent
on \(V_i\) is obtained by applying
Lemma~\ref{lem:regular-singular-descent} at a tangential index higher by
\(2\)---so \(n\) levels of descent reach down to tangential order
\(i+2n\) at the bottom, requiring control of \(\partial_t^kZ\) for
\(k\le i+2(n+1)\). The cap \(2\ell-1\) is exactly the
tangential range of \eqref{eq:tangential-z-bound}.

The hypothesis at step \(n\) is that \eqref{eq:Vi-n-form} has already
been established at every earlier order \(n'<n\), for all tangential
indices \(i'\) with \(i'+2(n'+1)\le2\ell-1\); the base case \(n=0\) is
\eqref{eq:Vi-base}.

For the inductive step, it suffices to verify that
\(F_i\in C^{n,\sigma}_{\loc}\): granting this,
Lemma~\ref{lem:regular-singular-descent} applied to
\eqref{eq:reg-sing-Vi} at order \(n\) yields \eqref{eq:Vi-n-form}; the
lemma's hypothesis on the principal coefficient is met by the
\(C^{\ell-1,\sigma}\)-regularity of \(b\) noted after
\eqref{eq:Fi-explicit}. The inductive constraint
\(i+2(n+1)\le2\ell-1\) forces \(n\le\ell-2\), which is exactly the
regularity available in the \(y\)-coefficients of
\eqref{eq:Fi-explicit}; hence those coefficients admit the required
\(n\) \(y\)-derivatives.

Membership in \(C^{n,\sigma}_{\loc}\) requires control of every mixed
derivative \(\partial_y^q\partial_t^rF_i\) with \(q+r\le n\). Because
\(c_\theta\), \(h\), and \(p_0''\) are independent of \(t\),
\eqref{eq:Fi-explicit} gives \(\partial_t^rF_i=F_{i+r}\), so the
problem reduces to estimating \(\partial_y^qF_{i+r}\) with
\(q+r\le n\). When \(q=0\) no normal derivative is involved, and the
tangential factors \(\partial_t^kZ\) of \(F_{i+r}\) are bounded by
\eqref{eq:tangential-z-bound}, since \(k\le i+r+2\le2\ell-1\). When
\(q\ge1\), applying the product rule to \eqref{eq:Fi-explicit} (with
\(i\) replaced by \(i+r\)) expands \(\partial_y^qF_{i+r}\) into a sum
of terms, each consisting of a \(y\)-derivative of a coefficient (each
loss covered by \(n\le\ell-2\)) times products of factors of the form
\(\partial_y^{q_a}\partial_t^{j_a}Z\) with \(q_a\le q\) and
\(j_a\le i+r+2\). If \(q_a=0\), such a factor is again controlled by
\eqref{eq:tangential-z-bound}. If \(q_a\ge1\), it equals
\(\partial_y^{q_a-1}V_{j_a}\), and the inductive hypothesis at
\((i',n')=(j_a,q_a-1)\) places it in \(C^{0,\sigma}_{\loc}\): one has
\(q_a-1<n\), and
\[
        j_a+2q_a\le i+r+2+2q\le i+2(n+1)\le2\ell-1,
\]
where the second inequality uses \(r+2q\le2n\), itself a consequence
of \(q+r\le n\). The extreme case \((q,r)=(n,0)\) saturates these
bounds: the leading \(Z\)-factor is then
\(\partial_y^n\partial_t^{i+2}Z=\partial_y^{n-1}V_{i+2}\), placed in
\(C^{0,\sigma}_{\loc}\) by the hypothesis at \((i+2,n-1)\) under the
borderline constraint \((i+2)+2((n-1)+1)=i+2(n+1)\le2\ell-1\); the
mixed cases \(r>0\) involve strictly fewer than \(n\) normal
derivatives and so rely only on earlier steps of the same triangular
induction. Hence \(F_i\in C^{n,\sigma}_{\loc}\), and Lemma
\ref{lem:regular-singular-descent} yields \eqref{eq:Vi-n-form}.

Setting \(j:=n+1\) in the unweighted half and \(j:=n+2\) in the weighted
half, \eqref{eq:Vi-n-form} translates to: \(\partial_y^j\partial_t^iZ
\in C^{0,\sigma}_{\loc}\) whenever \(j\ge1\) and \(i+2j\le2\ell-1\), and
\(y\,\partial_y^j\partial_t^iZ\in C^{0,\sigma}_{\loc}\) whenever
\(j\ge2\) and \(i+2j-2\le2\ell-1\). The \(C^{\ell,\sigma}\) regularity
below uses only the regime \(j+i\le\ell\); we record three sub-cases:
\begin{equation}
\label{eq:zji-bulk}
        \partial_y^j\partial_t^iZ\in C^{0,\sigma}_{\loc}([0,r_0^2/4)\times I),
        \qquad 1\le j,\quad j+i\le\ell-1,
\end{equation}
and, with at least one tangential derivative,
\begin{equation}
\label{eq:zji-edge}
        \partial_y^j\partial_t^iZ\in C^{0,\sigma}_{\loc}([0,r_0^2/4)\times I),
        \qquad 1\le j,\quad i\ge1,\quad j+i\le\ell,
\end{equation}
and
\begin{equation}
\label{eq:zji-top}
        y\,\partial_y^j\partial_t^iZ\in C^{0,\sigma}_{\loc}([0,r_0^2/4)\times I),
        \qquad 1\le j,\quad j+i=\ell .
\end{equation}
Since \(Z\) is bounded above and below by positive constants, smooth
composition transfers the bounds of \eqref{eq:tangential-z-bound} and
\eqref{eq:zji-bulk}--\eqref{eq:zji-top} to
\(\gamma_y=Z^{-1/(\theta+1)}\) and to
\(\gamma_y^{-\theta}=Z^{c_\theta}\). Hence
\eqref{eq:tangential-z-bound} and \eqref{eq:zji-bulk} give
\[
        \gamma(y,t)-\gamma_L(t)
        =\int_0^y \gamma_y(s,t)\dd s
        \in C^{\ell,\sigma}_{\loc}([0,r_0^2/4)\times I),
\]
and adding \eqref{eq:zji-edge} gives
\[
        \partial_t\!\left(\gamma(y,t)-\gamma_L(t)\right)
        =\int_0^y \gamma_{yt}(s,t)\dd s
        \in C^{\ell,\sigma}_{\loc}([0,r_0^2/4)\times I).
\]
Indeed, after one \(y\)-derivative the terms to control are
\(\partial_y^j\partial_t^i\gamma_y\) with \(j+i\le\ell\) and \(i\ge1\), while
the terms without \(y\)-differentiation involve only pure tangential
derivatives of \(\gamma_y\) up to order \(\ell+1\le2\ell-1\).
For the pressure trace, \eqref{eq:mass-relation} gives
\(m(\gamma(y,t),t)^\theta=p_0(y)\gamma_y(y,t)^{-\theta}\).
Since \(p_0\in C^{\ell,\sigma}\), for derivatives of order at most
\(\ell-1\) this product is \(C^{0,\sigma}\) by
\eqref{eq:tangential-z-bound} and \eqref{eq:zji-bulk}. For derivatives
of order \(\ell\), the only terms not already controlled are those in
which all \(y\)-derivatives fall on \(\gamma_y^{-\theta}\); there
\(p_0(y)=y\,h(y)\) absorbs the factor
\(y\,\partial_y^j\partial_t^i(\gamma_y^{-\theta})\), \(j+i=\ell\),
bounded by \eqref{eq:zji-top}. Thus the pressure trace is
\(C^{\ell,\sigma}\). The proof at \(b\) is identical after replacing
\(y\) by \(b-y\).
\end{proof}

\end{document}
